\chapter{Rotational CMC equations and their first integral}\label{sec:rotational-cmc-equations}
The graph equation and conserved quantity are described in \cite{Maz02}, Section 2.1, p.~4. The rotational metric and curvature formulas are also in \cite{BBM13}, pp.~2--3. We use trace mean curvature with $A=D\nu$ and derive the signs and arclength equations directly.

\begin{proposition}[Geometry of an arclength surface of revolution]\label{prop:rotational-arclength-curvature}
Let $(r(s),z(s))$ be a smooth plane curve on an interval, with $r>0$ and $(r')^2+(z')^2=1$. Write $r'=\cos\psi$, $z'=\sin\psi$ for a smooth angle lift. The immersion
\[
 F(s,\theta)=(r(s)\cos\theta,r(s)\sin\theta,z(s))
\]
has metric $ds^2+r^2d\theta^2$. Relative to $\nu=(z'\cos\theta,z'\sin\theta,-r')$, its principal curvatures are $\psi'$ and $z'/r$. Thus its trace mean curvature is $H=\psi'+\sin\psi/r$.
\end{proposition}
\begin{proof}
The coordinate derivatives are orthogonal, with lengths one and $r$. The indicated normal is orthogonal to both and has length one. Differentiation gives $D_s\nu=\psi'F_s$ and $D_\theta\nu=(z'/r)F_\theta$. These are the two eigenvalue equations for the shape operator. A smooth angle lift exists on any interval because the tangent takes values in the circle; it may wind, so no fixed principal branch of the angle is required.
\end{proof}

\begin{proposition}[The rotational CMC system and local uniqueness]\label{prop:rotational-cmc-system}
A rotational arclength immersion has constant trace curvature $H$ precisely when
\[
 r'=\cos\psi,\qquad z'=\sin\psi,\qquad
 \psi'=H-\frac{\sin\psi}{r}.
\]
For any initial data with $r>0$, this system has a unique smooth local solution. Its first integral is
\[
 c=r\sin\psi-\frac H2r^2.
\]
\end{proposition}
\begin{proof}
The curvature formula gives the system. Its right side is smooth for $r>0$, so the ordinary local existence and uniqueness theorem applies. Differentiation along a solution yields
\[
 c'=\cos\psi\sin\psi+r\cos\psi\left(H-\frac{\sin\psi}{r}\right)-Hr\cos\psi=0.
\]
Thus $c$ is constant throughout the interval of existence. The ODE theorem is a background analytic input; this statement gives no continuation through $r=0$.
\end{proof}

\begin{proposition}[The cylindrical graph equation]\label{prop:cmc-cylindrical-graph-ode}
For a positive smooth function $r=r(z)$, the cylindrical graph $(r(z)\cos\theta,r(z)\sin\theta,z)$, oriented by its outward radial normal, has trace mean curvature
\[
 H=\frac1{r\sqrt{1+r_z^2}}-\frac{r_{zz}}{(1+r_z^2)^{3/2}}.
\]
For constant $H$ its equation and first integral are
\[
 r_{zz}=\frac{1+r_z^2}{r}-H(1+r_z^2)^{3/2},\qquad
 c=\frac r{\sqrt{1+r_z^2}}-\frac H2r^2.
\]
\end{proposition}
\begin{proof}
Reparametrize the meridian by arclength increasing with $z$. Then $z_s=(1+r_z^2)^{-1/2}$, $r_s=r_z(1+r_z^2)^{-1/2}$, and the tangent angle satisfies $\psi_s=-r_{zz}(1+r_z^2)^{-3/2}$. Substitute into Proposition \ref{prop:rotational-arclength-curvature} and solve for $r_{zz}$. The conserved quantity is the preceding first integral with $\sin\psi=z_s$. At $H=1$, this is equation (1) in \cite{Maz02}, p.~4. The first integral used there is $-2c$.
\end{proof}

\begin{proposition}[The first integral from area minus volume]\label{prop:rotational-cmc-beltrami}
For fixed endpoint radii of a cylindrical graph, critical points of area minus $H$ times enclosed volume satisfy the ODE in Proposition \ref{prop:cmc-cylindrical-graph-ode}. The translation-invariant one-dimensional integrand is
\[
 L(r,p)=2\pi r\sqrt{1+p^2}-H\pi r^2,
\]
and its Beltrami quantity $L-pL_p$ equals $2\pi c$.
\end{proposition}
\begin{proof}
The lateral area element is $r\sqrt{1+r_z^2}\,d\theta\,dz$, and the volume between fixed axial planes is $\int\pi r^2dz$. Endpoint-fixed variations kill boundary terms. Direct differentiation gives
\[
 L_r=2\pi\sqrt{1+p^2}-2\pi Hr,\qquad L_p=\frac{2\pi rp}{\sqrt{1+p^2}}.
\]
The equation $(L_p)'=L_r$ simplifies to the stated graph ODE. Because $L$ has no explicit $z$ dependence, $(L-r_zL_p)'=r_z(L_r-(L_p)')=0$. Evaluating $L-r_zL_p$ gives $2\pi[r/\sqrt{1+r_z^2}-Hr^2/2]$. This is a stationarity calculation; it does not assert that a critical graph minimizes area under the volume constraint.
\end{proof}

\begin{lemma}[Factoring the positive-flux radius equation]\label{lem:unduloid-radius-factorization}
Fix $H>0$, $0<a<1/H$, $b=2/H-a$, and $c=Hab/2$. Then
\[
 r^2-\left(\frac H2r^2+c\right)^2
 =\frac{H^2}4(b^2-r^2)(r^2-a^2).
\]
In an arclength solution with this flux and $z'>0$, the radius belongs to $[a,b]$ and satisfies
\[
 (r')^2=\frac{H^2}{4r^2}(b^2-r^2)(r^2-a^2),\qquad
 z'=\frac H2r+\frac cr.
\]
\end{lemma}
\begin{proof}
Use $a+b=2/H$ and $ab=2c/H$ to compare coefficients in the polynomial identity. The conserved quantity gives $z'=Hr/2+c/r>0$, and $r'^2+z'^2=1$ gives the displayed radius equation. Its left side is nonnegative and $r>0$, so the factorization forces $a\le r\le b$. The squared equation alone is not a uniqueness statement at a turning radius; the smooth second-order system in Proposition \ref{prop:rotational-cmc-system} retains that information.
\end{proof}
